% =====================================================================
% Paper 2 -- Registered Report, STAGE 1 (introduction, methods, analysis plan).
% Written 2026-10-04 BEFORE any confirmatory result exists.
% Source of truth for every design element: PREREG_P2_discoverability.md (DRAFT, not frozen).
% Literature: research/lit_equation_discovery.md (codes S1..S65; each \cite key carries its
% S-code in paper2/references_p2.md).
% Number provenance: every number is quoted from the prereg or a results/test file named in a
% comment beside it. No confirmatory-grid result exists; none is implied.
% Aligned 2026-10-04 to the REVISED prereg (its s.10 "Pre-freeze revisions", items 1-20).
% =====================================================================
\documentclass[11pt]{article}
\usepackage[margin=2.5cm]{geometry}
\usepackage[T1]{fontenc}
\usepackage[utf8]{inputenc}
\usepackage{lmodern}
\usepackage{amsmath,amssymb}
\usepackage{booktabs}
\usepackage{array}
\usepackage{tabularx}
\usepackage[numbers,sort&compress]{natbib}
\usepackage{xcolor}
\usepackage[colorlinks=true,linkcolor=blue!50!black,citecolor=blue!50!black,urlcolor=blue!50!black]{hyperref}

\newcommand{\Da}{\mathrm{Da}}
\newcommand{\Pe}{\mathrm{Pe}}
\newcommand{\qs}{q^{*}}
\newcommand{\qsm}{q^{*}_{\mathrm{meas}}}
\newcommand{\verdict}[1]{``\textit{#1}''}

\title{When can AI discover adsorption kinetics? A discoverability map\\[4pt]
\large Registered Report --- Stage 1: introduction, methods and analysis plan}
\author{Abdulhamid Batayhi\\[2pt]
\small Department of Mechatronics Engineering, Faculty of Mechanical Engineering,\\
\small Y\i{}ld\i{}z Technical University, Istanbul, T\"urkiye\\
\small ORCID 0009-0007-0374-6755 \quad abdulbatayhi@gmail.com}
\date{Stage-1 draft, 4 October 2026}

\begin{document}
\maketitle

\begin{abstract}
Equation-discovery methods (sparse, weak-form, errors-in-variables and symbolic
regression, and universal differential equations) are proposed for learning the
adsorption rate law in fixed beds. The one published demonstration on breakthrough data
recovers the linear-driving-force law at a single kinetic rate and noise level, with the
isotherm supplied exactly. It is not known for which regimes such discovery succeeds, or
whether its failure coincides with the classical loss of practical identifiability of the
mass-transfer coefficient near local equilibrium; that loss is well documented and is not
claimed as new. We propose a pre-registered simulation study on data from a verified
non-isothermal column solver. It varies the Damk\"ohler number (nine levels,
$10^{-1}$--$10^{3}$), the P\'eclet number, noise, isotherm error, the observation model
(interior profiles with or without the solid loading, or outlet only), isotherm shape
(Langmuir or a MOF-303-like step) and the true law (constant $k$ or loading-dependent
$k(q)$). Discovery methods, each in its strongest configuration with documented inclusion
gates, are compared on one grid with profile-likelihood identifiability of $k$ given the
true law. Three confirmatory hypotheses, each with a falsifier and pre-declared verdict
words, ask whether the driving-force-to-error ratio collapses the boundary (H2a), whether
discovery fails before, with or after identifiability (H2b), and whether
errors-in-variables methods move the boundary (H2c); H2d, on the MOF case, is exploratory.
The boundary estimator was calibrated before any discovery run, giving a minimum
detectable effect of 0.25 decade. Pilots, reported separately, changed only the
generator, the harness and the feasible method set.
\end{abstract}

% ---------------------------------------------------------------------
\section{Introduction}
\label{sec:intro}

\subsection{The question}
A fixed-bed adsorption experiment reports concentrations and temperatures in time,
usually at the column outlet. The quantity an engineer needs from it is the uptake rate
law, most often written as the linear driving force (LDF), $\partial q/\partial t =
k\,(\qs(c,T) - q)$. AI equation-discovery methods offer to find that law, not only its
coefficient. Sparse regression over a library of candidate terms
\citep{brunton2016,rudy2016}, its weak-form \citep{messenger2021ode,messenger2021pde,
reinbold2020}, ensemble \citep{fasel2022}, Bayesian \citep{hirsh2021}, implicit
\citep{kaheman2020} and exact mixed-integer \citep{bertsimas2022} variants, symbolic
regression \citep{cranmer2023,udrescu2020} and neural surrogates embedded in the
governing equations \citep{santana2023} all make that offer. This paper asks a narrower
question, which can be tested:

\begin{quote}
\emph{For which regimes can the adsorption rate law be discovered from data a column
experiment can produce, and does discoverability fail at the same boundary as the
classical practical identifiability of the rate constant, or earlier?}
\end{quote}

\subsection{What is known and what is not new}
\label{sec:notnew}
The physical core of the expected failure is classical and is not claimed here. Near
local equilibrium the mass-transfer coefficient is poorly determined by column data.
Valocchi derived local-equilibrium validity criteria from temporal moments, in which
velocity, dispersion and rate jointly decide whether local equilibrium holds
\citep{valocchi1985}. Hansen and Vesselinov showed that those criteria really test
whether mass-transfer-driven dispersion is negligible next to hydrodynamic dispersion
\citep{hansen2017}. Katsoulas et al.\ derived asymptotically that, near local
equilibrium, finite mass-transfer rates enter only through an apparent dispersion
coefficient \citep{katsoulas2023}. In adsorption engineering, Knox et al.\ found that the
LDF coefficient and the axial dispersion can be extracted wrongly from outlet
breakthrough, and that interior measurements are needed \citep{knox2016}. Ward and Pini
attributed about 70\,\% of breakthrough output variability to the isotherm parameters
\citep{ward2022}. Taylor et al.\ mapped global parameter sensitivity of sorbing transport
across P\'eclet--Damk\"ohler regimes and found that equilibrium sorption dominates
\citep{taylor2026}. The counter-evidence is also classical. In the constant-pattern
regime, a lumped rate coefficient \emph{is} extracted from frontal breakthrough
consistently with shock-layer theory \citep{miyabe2000}. ``$k$ is unidentifiable at high
Damk\"ohler number'' is therefore not true without conditions. Any boundary must be
conditioned on dispersion, isotherm shape, noise and isotherm error, and those are the
axes of this design.

The generic failure of sparse regression on such data also has precedents.
Single-trajectory identifiability of linear-in-parameter systems, the SINDy model class,
depends on the trajectory's geometry; estimation sensitivity is governed by a condition
number tied to the data's spatial confinement \citep{stanhope2014}. The smallest
eigenvalue of the data moment matrix caps recovery for both SINDy and PySR
\citep{gallo2026}. SINDy becomes ill-conditioned on slow manifolds, where the recoverable
object is the manifold itself \citep{delgado2025}. Sparse identification shows noise- and
sparsity-induced phase transitions \citep{klishin2024}. When trajectories cross time-scale
regimes, SINDy fails to find a single global model \citep{muhammed2025}. Stiffness and
timescale separation break rate-law discovery in chemical kinetics
\citep{prabhu2025,huang2021}. Collinear library functionals make the discovered model
non-unique \citep{delahunt2021}. One benchmark points the other way: on attractor-filling
chaotic data, SINDy performance showed no significant dependence on scale separation
\citep{kaptanoglu2023}. Breakthrough data differ from that setting. A breakthrough run is
a transient single trajectory confined near the slow manifold $q \approx \qs(c,T)$, the
setting of \citep{stanhope2014,gallo2026,delgado2025}. The design measures this
confinement directly through library conditioning (\S\ref{sec:mech}).

\subsection{Prior equation discovery for adsorption}
The closest prior work is Santana et al.\ \citep{santana2023}. A neural network replaced
the solid--fluid mass-transfer term of a fixed-bed model (a universal differential
equation, UDE), was trained on outlet breakthrough curves with 5\,\% Gaussian noise, and
was then read out by sparse regression (BIC-selected) and genetic-programming symbolic
regression. For the LDF ground truth the recovered form was $p_1 + p_2 q + p_3\qs$ with
$p_2 \approx -0.22$ and $p_3 \approx 0.22$. The isotherm was supplied exactly, and one
kinetic rate and one noise level were studied \citep{santana2023}. The same group used a
universal ODE on experimental multicomponent fixed-bed data \citep{nogueira2022} and
screened ten hybrid placements of neural networks in chromatography models
\citep{frandsen2025}. Differentiable finite-volume hybrids now learn constitutive laws
from transport experiments \citep{jessop2026}. On the equilibrium side, neural
surrogates infer the isotherm from column dynamics \citep{praditia2022,subraveti2022}.
Bootstrap confidence intervals then reveal distinct learned sorption closures with similar
predictions, which is practical non-identifiability of a learned closure
\citep{scheurer2026}. Symbolic regression has rediscovered isotherms
\citep{haghpanah2023,sharlin2024,cornelio2021,fox2023}. Kinetic rate-law discovery by
symbolic regression with information criteria and design of experiments exists for batch
catalysis \citep{servia2024,servia2025}. Function-level identifiability of UDEs has been
separated from parametric identifiability \citep{loman2025}, and statistical tests have
been used to decide whether isotherm parameters in liquid chromatography must be
state-dependent \citep{katsoulas2026}. Three fetched database searches made while compiling
our literature review (two OpenAlex queries and one arXiv query) found no
adsorption-\emph{kinetics} discovery study other than \citep{santana2023}. This supports, but does not prove, that the
regime question is open.
% Source: research/lit_equation_discovery.md, "Negative-search evidence" 1 and 2, synthesis (a).

For metal--organic framework (MOF) water sorption the gap is sharper. Water uptake on
MOF-303 has been fitted with an LDF model by hand, and the specific diffusion rate falls
markedly where the isotherm inflects \citep{abuelmaaty2024}. A loading-dependent $k(q)$ at
an isotherm step is an untested discovery target. A 2026 review of AI for MOF water
harvesting does not mention kinetics discovery in its abstract \citep{coyle2026}.

\subsection{What this study adds}
Three things are new, if the hypotheses survive. The first is a quantitative
discoverability boundary for the strongest available discovery methods, with a mechanism:
the realised driving-force fraction against the combined error of the observations and of
the isotherm the method is given. The second is a like-for-like comparison of that
boundary with profile-likelihood identifiability \citep{raue2009} on the same grid, with
the same noise, isotherm error and observation model. The third is a test of whether an
errors-in-variables estimator moves the boundary. Isotherm error is an error in a
regressor, not in the target, and the paper treats it so: the grid uses classical mixed
least-squares--total-least-squares fitting \citep{golub1980,golub1987} inside the exact
best-subset search. Errors-in-variables equation-discovery methods
\citep{fung2026,bortz2023} are reported on their own published examples. A negative result
on any of these is publishable in the registered-report format, and the design is built to
make it informative. Every null carries its minimum detectable effect, and every method
that cannot be run in its strongest configuration is reported as such.

% ---------------------------------------------------------------------
\section{Hypotheses}
\label{sec:hyp}
Let $\delta$ be the realised driving-force fraction, computed from ground truth as the
maximum over the observed $(z,t)$,
\begin{equation}
\delta = \max_{(z,t)\ \mathrm{observed}}\,\lvert \qs(c,T) - q\rvert / q_{\max}.
\end{equation}
Let $\varepsilon_{\mathrm{eff}}$ be the combined relative error of the equilibrium the
method is given (isotherm error $\varepsilon$), of the observations (noise $\sigma$) and,
for O2, of the reconstructed state:
\begin{equation}
\varepsilon_{\mathrm{eff}} = \sqrt{\sigma^{2} + \varepsilon^{2} + \varepsilon_{\mathrm{rec}}^{2}},
\end{equation}
where $\varepsilon_{\mathrm{rec}} = 0$ for O1. For O2 it is the RMS of $(\hat q -
q)/q_{\max}$ of the mass-balance-inverted uptake (M1b, \S\ref{sec:compet}) against ground
truth. It is computed once per cell from the noise-free ($\sigma = 0$) O2 observation, so
it measures the inversion's structural error; the noise is already counted in $\sigma$.
Then $R = \delta/\varepsilon_{\mathrm{eff}}$, computed per cell.
% PREREG_P2 s.4.2 "Combined error and driving force (for H2a)"; s.4.3 "eps_rec".

\paragraph{H2a: the boundary is a driving-force-versus-error boundary.}
The probability of successful discovery, $P(\text{success})$, collapses onto one curve
when plotted against $R$, across $\Da$, $\Pe$, isotherm shape and observation model (O1
and O2; O3 has no discovery method beyond M1's demonstration). The
50\,\% crossing lies at $R^{*}$ with a spread of \emph{less than half a decade} across
conditions.
\emph{Falsifier:} an $R^{*}$ spread of one decade or more across conditions.
\emph{Pre-declared words:} \verdict{R collapses the boundary}, or
\verdict{R does not collapse the boundary; the boundary depends additionally on
\{named axis\}}.
A spread of at least half a decade but less than one decade neither confirms nor
falsifies H2a and is reported in the pre-declared words \verdict{inconclusive: R narrows
the boundary to a spread of \{x\} decades}.
% PREREG_P2 s.4.4 "Verdict words for the gaps between thresholds".

\paragraph{H2b: discoverability versus identifiability.}
For each slice ($\sigma$, $\varepsilon$, observation model, isotherm), we compare
$\Da^{*}_{\mathrm{disc}}$, the 50\,\% crossing of $P(\text{success})$ for the best method,
with $\Da^{*}_{\mathrm{ident}}$, the boundary beyond which $k$ is no longer practically
identifiable given the \emph{true} law. Both are 50\,\% crossings from the same logistic
estimator (\S\ref{sec:ident}). \textbf{H2b is confirmatory on O1 only.} There both
boundaries rest on 20 replicates per level, and the 0.25-decade MDE applies. On O2 and O3,
identifiability has 3 replicates on a subgrid, so those comparisons are reported
descriptively, each with its own bootstrap interval, and carry no verdict word. On O1 there
are three pre-declared outcomes, reported in these words:
\begin{itemize}
\item \verdict{discovery fails before identifiability}: $\Da^{*}_{\mathrm{disc}} <
\Da^{*}_{\mathrm{ident}}$ by at least the MDE;
\item \verdict{discovery fails with identifiability}: the absolute difference is less than
the MDE;
\item \verdict{discovery outlives identifiability}: $\Da^{*}_{\mathrm{disc}} >
\Da^{*}_{\mathrm{ident}}$ by at least the MDE. The method then finds the law's form where
its coefficient is undetermined.
\end{itemize}
When either boundary has no crossing inside the tested range, no verdict word is issued and
the range tested is reported instead (\S\ref{sec:analysis}).

\paragraph{H2c: errors-in-variables methods move the boundary.}
The errors-in-variables (EIV) method on the grid (M5, \S\ref{sec:compet}) shifts
$\Da^{*}_{\mathrm{disc}}$ by at least a factor of 2 relative to the best non-EIV method at
$\varepsilon > 0$. The comparison is like for like: both methods in the strong form (M5 is
strong-form only), on O1. The best non-EIV method over both forms is reported alongside,
descriptively.
\emph{Falsifier:} a shift below the MDE at every $\varepsilon > 0$.
\emph{Pre-declared words:} \verdict{EIV moves the boundary by \{factor\}}, or
\verdict{EIV does not move the boundary detectably (MDE \{x\})}. A shift at or above the
MDE but below 0.30 decade (a factor of 2) is reported as \verdict{EIV moves the boundary
by \{factor\}, detectably but by less than the predicted factor 2}.
% PREREG_P2 s.4.4 "Verdict words for the gaps between thresholds".

\paragraph{H2d: the MOF case (exploratory, and labelled so).}
For a step (Type V) isotherm with a rate that drops at the step, as reported for MOF-303
\citep{abuelmaaty2024}, is the state dependence $k(q)$ discoverable, and in which
($\sigma$, observation) cells? No directional prediction is made, and no confirmatory claim
rests on H2d. Its success criterion is fixed in \S\ref{sec:metrics}.

% ---------------------------------------------------------------------
\section{Methods}
\label{sec:methods}

\subsection{Ground truth}
\label{sec:truth}
Ground truth comes from the verified finite-difference reference solver of the companion
paper \citep{paper1}. It is a one-dimensional, axially dispersed, non-isothermal column
model, and its implementation is verified against closed-form solutions (the companion
paper's L0 checks). The uptake law is one of two:
\begin{align}
\text{L1:}\quad & \frac{\partial q}{\partial t} = k\,\bigl(\qs(c,T) - q\bigr),\\
\text{L2:}\quad & \frac{\partial q}{\partial t} = k(q)\,\bigl(\qs(c,T) - q\bigr),\qquad
k(q) = k_0\Bigl[1 - 0.8\,s\!\Bigl(\frac{q/q_{\max}-\theta}{0.05}\Bigr)\Bigr],
\end{align}
where $s$ is the logistic function and $\theta$ is the fractional loading at the
isotherm's step midpoint, computed from the isotherm at the feed temperature and recorded
per solve. Under L2 the rate falls to 20\,\% of $k_0$ across the step. L2 is implemented
through the solver's state-dependent-rate hook, whose default path is verified unchanged.
% PREREG_P2 s.4.1.
Non-isothermal physics as in the companion paper is primary; an isothermal control is
secondary. There are two isotherms. ``Langmuir'' is the companion paper's default
favourable isotherm. ``Step'' uses the cited MOF-303 step position and capacity, with
column geometry and flow held at the default so that only the isotherm changes. Ground
truth is stored with a manifest that records every parameter, the stoichiometric time
$t_{\mathrm{stoich}}$, the horizon, and a per-solve breakthrough check
(\S\ref{sec:pilots}).

\subsection{Design axes}
\label{sec:axes}
Table~\ref{tab:axes} lists the axes. The primary Damk\"ohler number is
$\Da = k\,t_{\mathrm{stoich}}$, the number of kinetic time constants per front passage. The
residence-time form $kL/u$ is reported alongside but is not the axis. In this column the
residence time is 0.4\,s, against a front passage of about 7200\,s, so $kL/u = 0.4\,
\mathrm{s}\times k$. That is $4\times10^{-5}$ to 4 over the pilot's $k = 10^{-4}$--$10$\,s$^{-1}$,
and at most about 0.04 over the targeted kinetic range, so it cannot separate the regimes
this study is about (\S\ref{sec:pilots}).
% PREREG_P2 s.4.2; results/p2_pilot_generator.json: Da_res = 0.4 s x k (4e-5 ... 4);
% t_final 10800.6 s = 1.5 t_stoich -> t_stoich ~ 7200 s.
\begin{table}[t]
\centering
\caption{Design axes (pre-registration \S4.2).}
\label{tab:axes}
\begin{tabularx}{\textwidth}{>{\raggedright\arraybackslash}p{4.2cm}X}
\toprule
axis & levels \\
\midrule
$\Da = k\,t_{\mathrm{stoich}}$ & 9, log-spaced, $10^{-1}$--$10^{3}$ (half-decade spacing) \\
$\Pe$ (via the axial-dispersion multiplier) & 3: $\times 0.2$, $\times 1$, $\times 5$ \\
noise $\sigma$, i.i.d.\ Gaussian, relative to channel range & 0, 0.5\,\%, 2\,\%, 5\,\% \\
isotherm error $\varepsilon$, smooth multiplicative perturbation of $\qs$ & 0, 0.5\,\%, 2\,\%, 5\,\% \\
observation model & O1 interior $c,q,T$; O2 interior $c,T$ ($q$ latent); O3 outlet $c(t),T(t)$ only \\
isotherm & favourable Langmuir; MOF-303-like step \\
true law & L1 constant $k$; L2 $k(q)$ (step isotherm only) \\
\bottomrule
\end{tabularx}
\end{table}

Noise and isotherm error are applied after the solve. Ground truth therefore needs only
the ($\Da$, $\Pe$, isotherm, law) solves: 9 $\times$ 3 $\times$ 3 = 81, over three
(isotherm, law) pairs (Langmuir--L1, step--L1, step--L2). These 81 solves are done, at 3--7\,s
each at the grid's $k$ (the generator pilot measured up to about 300\,s at $k \ge
0.03$\,s$^{-1}$). Every solve reached 95\,\% breakthrough, some after the horizon was
extended (\S\ref{sec:pilots}). This is ground truth, not a confirmatory result: no
discovery method has seen it.
% PREREG_P2 s.4.2; p2_generate.log "done: 81 cells; without 95 % breakthrough: none".
Replicates are 20 independent ($\sigma$, $\varepsilon$) realisations per cell for M2--M6
and M9 (M8 is not run on the grid; \S\ref{sec:compet}). M6 runs on a separate isothermal
control: 9 solves, the 9 $\Da$ levels $\times$ $\Pe \times 1$ $\times$ Langmuir, with the
energy balance switched off (T held at the feed) and the same horizon rule. These solves
are done; all reached 95\,\% breakthrough. M7, if its inclusion gate passes, needs a Julia search per replicate and so
runs with 3 replicates on a subgrid: all $\Da$ $\times$ $\sigma \in \{0.5\,\%, 2\,\%\}$
$\times$ $\varepsilon \in \{0, 2\,\%\}$ $\times$ both isotherms $\times$ $\Pe \times 1$
$\times$ O1, the only observation model M7 can see (Table~\ref{tab:obs}). M1 runs only
as the minimal demonstration declared in \S\ref{sec:compet}, nothing more.

\paragraph{Observation model, exactly (\texttt{p2\_observe.py}).}
All methods and the identifiability analysis read data only through one observation
function, so every method sees identical corruption. O1 and O2 observe 20 equally spaced
interior probes, and O3 observes the outlet; every model samples 200 times uniform over
the solved horizon. Noise is i.i.d.\ Gaussian with standard deviation $\sigma$ times the
channel's range in that cell's ground truth. The measured isotherm is $\qsm(c,T) =
\qs(c,T)\,[1 + \varepsilon\, g(c/c_{\mathrm{in}})]$, with $g$ a smooth random function of
4 Fourier modes on $[0,1]$, normalised to unit RMS and drawn once per replicate. Every
draw is deterministic per (cell, replicate, $\sigma$, $\varepsilon$, observation model).
The true $\qs$ is never passed to a method.
% PREREG_P2 s.4.2 "Observation model, exactly".

\subsection{Candidate libraries}
\label{sec:libs}
The two libraries were fixed before the freeze to the implemented and
known-answer-tested harness.
\begin{itemize}
\item \textbf{Lib-B (equilibrium-informed, primary):} $\{1,\ c,\ q,\ T,\ c\,q,\ q^{2},\
c^{2},\ \qsm\}$. Including $\qsm$ is legitimate because it is the \emph{measured} isotherm,
carrying error $\varepsilon$. The earlier exploratory script that motivated this study
included the exact generating law, which was a defect. An earlier draft of the design
specified degree-3 polynomials, $\exp(-\theta/T)$ factors and $\qs q$ products; the tested
harness uses the smaller library above. A smaller library makes discovery easier, so the
boundaries measured here are \emph{optimistic with respect to library size}.
\item \textbf{Lib-A (agnostic):} Lib-B without $\qsm$. Under L1 the true law is not
representable in Lib-A, because $\qs$ is not polynomial, so Lib-A cannot succeed by the
primary metric. It is reported as held-out fidelity only. The agnostic route to the law's
form is M6 (SINDy-PI) on the isothermal control, which finds the implicit Langmuir--LDF
structure without any isotherm.
\end{itemize}
% PREREG_P2 s.3 (libraries); p2/features.py lib_b.

\subsection{Methods compared, each in its strongest configuration}
\label{sec:compet}
The rule is that each method is run in its strongest configuration. A method that cannot be
so run is reported as such and is never replaced silently by a weaker one. Implementation
choices were fixed against known-answer systems before any grid cell was touched.

\begin{description}
\item[M1, the S1 pipeline reproduced.] A UDE with a network in the uptake term, trained
through the column PDE and read out by sparse and symbolic regression
\citep{santana2023}. Its settings are at least those of \citep{santana2023}, and stronger
where those are weak. Feasibility was resolved by a pilot (\S\ref{sec:pilots}): M1 is
\emph{not run at the declared scale}. It runs only as a labelled demonstration on a
minimal subgrid: Langmuir, $\Pe \times 1$, $\sigma = 2\,\%$, $\varepsilon = 0$, O3, the
9 $\Da$ levels, one replicate, L-BFGS, at most 50 steps. The O3 comparison rests on the
profile likelihood. This is a compute limit, not a finding about UDEs; on a GPU the same
code is expected to be feasible.

\item[M1b, mass-balance inversion (O2).] From interior $c(z,t)$ the uptake rate follows
from the gas-phase balance,
\begin{equation}
\frac{\partial q}{\partial t} = -\frac{\epsilon}{(1-\epsilon)\rho_p}\Bigl(
\frac{\partial c}{\partial t} + u\frac{\partial c}{\partial z} - D_L
\frac{\partial^2 c}{\partial z^2}\Bigr),
\end{equation}
with $q$ recovered by time integration from $q(z,0) = 0$. This is how an experimentalist
with interior probes obtains uptake. It needs no training, and every O1 method then runs
on its output, which makes it the strongest O2 competitor. Its coefficients come from the
solver's own definition, so the inversion cannot drift from the forward model. Time
derivatives use the same Savitzky--Golay rule as O1. Space derivatives use centred
finite differences, because the probe spacing is fixed by the observation model and no
smoothing window in $z$ may be tuned to the answer.
% PREREG_P2 s.3 M1b; p2/massbal.py.
On a fresh solve that is not a grid cell, the inversion recovers $q$ to within the 3\,\%
of $q_{\max}$ bound recorded in the pre-registration. The test measures a maximum interior
error of 2.4\,\% (\S\ref{sec:verified}).
% tests/test_p2.py::test_mass_balance_inversion_recovers_uptake_from_gas_data (comment:
% "measured 0.024"); commit d142180.

\item[M2, weak-form SINDy \citep{messenger2021ode,messenger2021pde}; M2b, its strong-form
counterpart.] The weak transform uses compactly supported test functions $\phi(t) =
((t-a)(b-t))^{8}$ and integration by parts, so the target needs no derivative of noisy
data. The test-function support is swept over $\{31, 61, 121\}$ samples and chosen per
cell by the held-out weak-system residual on every fifth probe. The law is then refitted
on all probes. A support is admissible only if its training system has at least 3
equations per library term. On the known-answer test, a long support left an
under-determined weak system that fitted a wrong, dense law with zero residual, which a
residual criterion would otherwise prefer. M2b is the derivative-based regression on a
Savitzky--Golay $\mathrm{d}q/\mathrm{d}t$. Both forms run every sparse solver (STLSQ, M3,
M4).
% PREREG_P2 s.3 M2 row; p2/run_o1.py choose_support (WEAK_SUPPORTS, MIN_ROWS_PER_TERM = 3);
% p2/weak.py; p2_grid_o1.py FORMS = ("strong", "weak").

\item[M3, ensemble or Bayesian SINDy \citep{fasel2022,hirsh2021}.] The inclusion
probability of each term is reported. The implemented configuration is bagged
sequentially thresholded least squares with an inclusion threshold of 0.6.
% PREREG_P2 s.3 Phase A; p2/methods.py ensemble(inclusion=0.6).

\item[M4, exact best-subset selection by enumeration.] This is the MIOSR objective
\citep{bertsimas2022}, solved exactly because the library is small and so needing no
commercial solver. The search is exhaustive over sparsity $\le 4$, BIC selects, and the
fit is \emph{unconstrained}. An earlier draft added sign constraints ($k > 0$, zero rate at
$q = \qs$). They were dropped before the freeze because ``zero rate at $q = \qs$''
imposes the very coefficient ratio $-1$ that the success metric tests, so M4 would
succeed partly by construction. Signs are judged by the metric, as for every method. BIC
carries a residual-variance floor at $2\times10^{-3}\times\mathrm{std}(\mathrm{d}q/
\mathrm{d}t)$, tied to the derivative estimator's accuracy at the grid's sampling (RMS
$8\times10^{-4}$ at 200 times). At zero noise, plain BIC kept terms of relative size
$10^{-4}$ by fitting the estimator's own systematic error. The floor was first set at
$10^{-3}$ from 400-sample tests, where M6 then bought a spurious $q^{2}$ term at zero noise
on 200 samples. $2\times10^{-3}$ is the smallest of $\{1,2,5\}\times10^{-3}$ that passes,
calibrated on the known-answer system only. These design facts were forced by
known-answer tests and are declared before the freeze.
% PREREG_P2 s.3 Phase-A design fact (i); p2/methods.py best_subset(floor=2e-3).

\item[M5, errors-in-variables.] \emph{On the grid, M5 is EIV best-subset.} It runs M4's
exact enumeration, but fits each support by classical mixed least-squares--total-least-squares
\citep{golub1987}, with the total-least-squares step computed by singular value
decomposition \citep{golub1980}. Exact columns (the constant) are fitted by least squares.
Error-carrying columns and the target are fitted by total least squares, after each is
scaled by its \emph{declared} error standard deviation. For the observed channels this is
$\sigma$ times the channel range. For $\qsm$ it is $\varepsilon\,\qs$ plus its
$c$-sensitivity. Products get first-order propagation. The target's standard deviation
comes from the Savitzky--Golay filter's exact noise gain. Selection uses M4's floored BIC.
M5 runs on the full O1 grid, with 20 replicates, in the strong form. Two known-answer tests
check it: total least squares removes the attenuation bias that ordinary least squares
suffers, and at zero error EIV best-subset recovers L1 with M4's support.
% PREREG_P2 s.3 M5 row and s.10; p2/eiv.py; p2_grid_extra.py m5.

\emph{ODR-BINDy \citep{fung2026} and WENDy-IRLS \citep{bortz2023} are reported on their
own published examples only.} Both learn \emph{autonomous} ODE systems, in which every
state is denoised and given its own equation. The target here is one rate equation with
exogenous $c$ and $T$ inside a PDE, and the dominant error is in a regressor, the measured
isotherm. Extending either method to that problem would be new method development. This
was decided before the freeze. The ODR-BINDy Python port, on the authors' Lorenz example
(500 points, 20\,\% noise), recovers the exact 7-term support. Its relative coefficient
error is $5.0\times10^{-3}$, against about $2.2\times10^{-3}$ read from the paper's heatmap,
and it is about 26 times slower (2385\,s). Before the freeze, five noise seeds separate
single-run scatter from a real gap. The statistic is the mean relative coefficient error
over the seeds, compared against $10^{-2.65}$. If the gap persists, the port is reported as
\verdict{port, accuracy below the reference}, with the gap stated. WENDy is reported as
\verdict{NOT RUN} if no implementation is verified on its authors' own example before the
freeze.
% PREREG_P2 s.3 ODR-BINDy port paragraph; results/odr_bindy_verification.json:
% rel_frobenius_error 0.00498, reported 0.00224, runtime_s 2384.75; p2_odr_seeds.py exists,
% results/odr_bindy_seeds.json does not yet exist (2026-10-04).

\item[M6, SINDy-PI \citep{kaheman2020} for rational laws.] M6 is run with Lib-A and
\emph{no} measured isotherm, on the isothermal control only (\S\ref{sec:axes}), because
$\exp(-\Delta H/RT)$ is not polynomial. Success is the implicit Langmuir--LDF structure
$\{\mathrm{d}q,\ c\,\mathrm{d}q,\ c,\ q,\ c\,q\}$ with correct signs. Its library is
$\{1,\ c,\ q,\ c\,q,\ q^{2},\ c^{2}\}$, which is Lib-A without $T$, because $T$ is constant
on the isothermal control. It runs at every $\sigma$ and at $\varepsilon = 0$ only, because M6
is given no isotherm and isotherm error does not apply. M6 is verified on
the known-answer system, with every coefficient (hence $b$ and $k$) recovered within
5\,\%.
% PREREG_P2 s.3 M6 row and s.10 item 10; tests/test_p2.py line 259:
% abs(coefs[key] - true) <= 0.05 * abs(true).

\item[M7, constrained symbolic regression with PySR \citep{cranmer2023}.] Julia is
installed only in an isolated environment. \emph{Inclusion gate, written before any M7
search:} PySR (binary $+,-,\times$; maximum size 12; 40 iterations; PySR's own ``best''
model selection) is run on the known-answer LDF system ($k = 0.02$; 0.5\,\% noise on $q$
\emph{and} on the measured $\qsm$) for 3 seeds. M7 runs on its subgrid if and only if success holds in at least 2 of 3 seeds,
\emph{and} every search finishes within 30\,min, \emph{and} peak memory stays under
3\,GB. Otherwise it is reported as NOT RUN, with the pilot's numbers, as ``infeasible on
this machine'' or ``failing its known-answer test'', in those words. The gate pilot is
queued behind the companion paper's runs.
% PREREG_P2 s.3 M7 row; p2_pilot_m7.py; results/p2_pilot_m7.json does not yet exist.

\item[M8, KAN symbolic extraction \citep{liu2024}.] M8 follows the original KAN
procedure. It is not a reimplementation of KANDy \citep{slote2026}, whose code was not
opened. It is a secondary method. \emph{It fails its known-answer test at zero noise in
every configuration tried.} The polynomial library at three sparsity penalties routes the
law through $q^{2}$. Raw variables on one trajectory use $q$ alone. Raw variables on three
trajectories return $\{1, \qs, {\qs}^{2}\}$. This is consistent with the reported
fragility of KANs on multiplicative terms \citep{spotorno2026}. M8's grid run was gated on
cost. Its best configuration is raw variables with sparsity penalty $10^{-3}$. Its subgrid
is all $\Da$ $\times$ $\sigma \in \{0.5\,\%, 2\,\%\}$ $\times$ $\varepsilon \in \{0, 2\,\%\}$
$\times$ both isotherms $\times$ $\Pe \times 1$ $\times$ O1, with 3 replicates, which is 216
fits. The rule, written into the cost pilot before it ran, was that M8 runs only if one fit
at the grid's size takes at most 300\,s, so at most 18\,h for the subgrid. \emph{The pilot
measured 1144\,s, so M8 is NOT RUN on the grid.} Over the subgrid that would have been
about 69\,h. This is a machine limit, reported in those
words. It was measured while the machine was shared with other long jobs, which the paper
states. On the same known-answer data the timed fit returned only a constant, consistent
with the documented failure. M8 is reported on its known-answer tests only.
% PREREG_P2 s.10 "M8 cost gate"; results/p2_cost_gates.json: sec_per_fit 1144.4,
% limit_s 300, fits_on_subgrid 216, projected_h 68.7, support_on_known_answer ["1"];
% tests/test_p2.py::test_kan_symbolic_recovers_ldf_at_zero_noise (xfail, strict=True).

\item[M9, slow-manifold discovery \citep{delgado2025}.] M9 regresses the \emph{state} $q$,
not its rate, on the $q$-free features $\{1,\ c,\ T,\ c^{2},\ \qsm\}$ by exact best-subset
selection. It is expected to win at
high $\Da$ by returning the isotherm (plus apparent dispersion). It is reported as
\emph{what is discoverable there}, not as rate-law success. Its own success criterion is
given in \S\ref{sec:metrics}.
\end{description}

\paragraph{Which method can see which observation model.}
This was decided before the freeze; a method is never run where its input does not exist
(Table~\ref{tab:obs}). O1 is the optimistic ceiling. O2 and O3 are where the
experimentally realistic answer lies, and O3 is reported first.

\begin{table}[t]
\centering
\caption{Method--observation applicability (pre-registration \S3).}
\label{tab:obs}
\small
\begin{tabular}{>{\raggedright\arraybackslash}p{5.6cm}ccc}
\toprule
 & O1 ($c,q,T$ interior) & O2 ($c,T$ interior) & O3 (outlet only)\\
\midrule
M2, M2b, M3, M4, M6, M8$^{a}$ (regression on $q$, $\mathrm{d}q/\mathrm{d}t$) & yes & no & no\\
M5 EIV best-subset (mixed LS--TLS) & yes & no & no\\
M7 PySR & yes & no & no\\
M9 slow manifold & yes & no & no\\
M1b inversion, then every O1 method & --- & yes & no\\
M1 UDE + sparse/symbolic readout & yes & yes & yes\\
profile likelihood (law known) & ODE per probe & PDE solve & PDE solve\\
\bottomrule
\multicolumn{4}{l}{\footnotesize $^{a}$ Applicable, but NOT RUN on the grid: it failed its cost gate (\S\ref{sec:compet}).}
\end{tabular}
\end{table}

\subsection{Success metric and secondary accuracy}
\label{sec:metrics}
\textbf{Success (primary)} means that the selected support equals the true law's support
\emph{and} every coefficient has the correct sign. For L1 in Lib-B the support must be
exactly $\{\qsm, q\}$, with coefficient ratio $-1$ within 10\,\%.
\textbf{Accuracy (secondary)} is $\lvert\hat k - k\rvert/k \le 10\,\%$, with $\hat k$ the
negative of the $q$ coefficient. $P(\text{success})$ is computed per cell over
replicates. $\Da^{*}_{\mathrm{disc}}$ is obtained from a logistic fit in $\log_{10}\Da$,
with a bootstrap confidence interval over replicates (\S\ref{sec:analysis}). A method that
raises an error in a cell is recorded as a failure with the error message, never dropped.
% p2/metric.py success_L1, k_accuracy; p2_grid_o1.py records errors as failures.

\textbf{Success for the other targets}, fixed before the freeze:
\begin{itemize}
\item \emph{M9:} support $\{\qsm\}$ with coefficient 1 within 5\,\%.
\item \emph{L2 (H2d, exploratory), operationally:} with $\qsm$ held fixed, the fitted law
$f$'s implied rate $\hat k(q) = -\partial f/\partial q$ is evaluated at the observed samples
and averaged in 10 equal bins of $q/q_{\max}$ over the observed range. The state dependence
is \emph{discovered} if and only if three conditions hold: the support contains $\qsm$ and
$q$; $\hat k$ decreases across the bins (Spearman $\rho < -0.8$); and $\hat k$ in the top
bin is at most 0.6 of $\hat k$ in the bottom bin. The true $k$ falls to 0.2 of its starting
value. This is reported as a rate per cell, with no boundary hypothesis.
\item \emph{Lib-A fidelity:} the relative RMS error of the predicted
$\mathrm{d}q/\mathrm{d}t$ (strong form) on the held-out probes (every fifth), per cell.
\end{itemize}
% PREREG_P2 s.4.3 ("L2 success, operationally", "Lib-A fidelity"); p2/metric.py
% success_manifold, success_L2 (n_bins=10, rho_max=-0.8, drop_max=0.6).

\textbf{The best method and its selection.} $\Da^{*}_{\mathrm{disc}}$ is estimated per
method per slice. H2a and H2b use the method with the largest $\Da^{*}_{\mathrm{disc}}$ in
the slice. Because that choice is made on the same replicates, its interval is a
bootstrap that resamples replicates \emph{jointly} across methods and re-selects the best
method inside every draw, so the interval carries the selection. This is the same
correction as the companion paper's selection-adjusted intervals. Every method's own
$\Da^{*}_{\mathrm{disc}}$ is reported alongside.
% PREREG_P2 s.4.3 "The 'best method' and its selection".

\subsection{Identifiability criterion}
\label{sec:ident}
Identifiability is assessed by the profile likelihood of $k$ \citep{raue2009} with the true
law known, at the same $\sigma$, $\varepsilon$ and observation model.
$k$ is \emph{practically identifiable} if and only if the 95\,\% profile interval
(likelihood-ratio threshold 3.84) lies within $[k/2, 2k]$. In O1 the profile integrates the
known-law ODE at each probe, driven by observed $c$ and $T$ (monotone cubic interpolation).
In O2 and O3 each profile point is a full column solve, scored against the interior $c$,
$T$ probes (O2) or the outlet $c(t)$, $T(t)$ (O3), on a coarse 25-point $k$ grid, with 3
replicates on the declared subgrid. Both PDE profiles are implemented and pass
known-answer tests.
% p2/identifiability.py profile_interval_probes (O2), profile_interval_outlet (O3);
% p2_grid_o23.py o2ident, o3. In both cases the profile refines adaptively between
the coarse neighbours of its interval. It makes up to 6 passes of 41 points (ODE profile)
or 21 points (PDE profile), each pass narrowing to the new neighbours. At low noise the
interval is narrower than one coarse step, and a single fixed refinement collapsed onto one
point below the truth on the known-answer test. The profile is given the \emph{measured}
isotherm, so identifiability faces the same isotherm error as discovery.
% PREREG_P2 s.3 Phase-A design fact (ii); p2/identifiability.py; p2_grid_o23.py run_o3:
% phys_obs.qstar_fn = obs["qstar_meas"] (information parity).

\textbf{$\Da^{*}_{\mathrm{ident}}$ is estimated exactly like $\Da^{*}_{\mathrm{disc}}$.}
Per replicate, $k$ is identifiable or not under the $[k/2, 2k]$ criterion.
$P(\text{identifiable})$ is fitted against $\log_{10}\Da$ by the same logistic estimator
(\S\ref{sec:analysis}), and $\Da^{*}_{\mathrm{ident}}$ is its 50\,\% crossing. The two
boundaries compared in H2b are therefore measured by one estimator.
% PREREG_P2 s.4.3 "Da*_ident, estimated exactly like Da*_disc".

\subsection{Mechanistic predictors}
\label{sec:mech}
For each cell we record $\delta$ from ground truth, the condition number of the library
Gram matrix, and its smallest eigenvalue, following
\citep{gallo2026,stanhope2014}. These test whether library conditioning, set by how
closely the data hug $q = \qs$, predicts the boundary. They do not enter any verdict.

\subsection{Implemented and verified components}
\label{sec:verified}
The harness is verified on known-answer systems by a test suite (\texttt{tests/test\_p2.py}).
None of these systems is a grid cell. The components and their checks are:
\begin{itemize}
\item a synthetic LDF series obeying its law to $10^{-3}$ of the rate scale, and a
Savitzky--Golay derivative exact to $10^{-3}$ on a smooth series;
\item the success metric, which rejects a spurious term of size $10^{-6}$, a wrong sign and
a ratio error of 50\,\%, and accepts a 7.5\,\% ratio error;
% test_success_requires_exact_support_signs_and_ratio: q = -0.03 vs 0.02 rejected;
% -0.0215 accepted.
\item M4 exact enumeration, STLSQ and the bagged ensemble, each recovering the LDF law at
zero noise (M4 with $k$ error below 2\,\%);
\item the weak form, recovering the law at zero noise and at 2\,\% noise with $k$ error
below 5\,\%, and beating the strong form at 2\,\% noise;
\item the O1 runner in both forms with all three sparse solvers;
\item M1b inversion to within 3\,\% of $q_{\max}$ on a fresh solve, and the O2 pipeline
recovering the law when fed the true $q$. Fed the inverted $q$ at that $\Da$, it does not
recover it, a pilot fact declared in \S\ref{sec:pilots};
\item the differentiable column used for M1, matching the reference solver to
$5\times10^{-3}$ on the exit curve;
\item the O1 profile likelihood (isothermal, and non-isothermal with temperature), the O2
interior-probe profile and the O3 outlet profile, each bracketing $k$ within the
pre-registered $[k/2, 2k]$ criterion in the kinetic regime;
\item EIV best-subset: total least squares removing the attenuation bias of ordinary least
squares, and recovery of L1 with M4's support at zero error;
\item the L2 success metric and the Lib-A fidelity metric;
\item SINDy-PI recovering the implicit Langmuir--LDF structure with no isotherm supplied,
every coefficient within 5\,\%;
\item the boundary estimator's coverage and bias (\S\ref{sec:analysis}), and the verdict
functions for H2a and H2b at the pre-declared thresholds;
\item M8's known-answer test, kept as an expected failure that will break the suite if a
library change makes it pass.
\end{itemize}

% ---------------------------------------------------------------------
\section{Analysis plan}
\label{sec:analysis}

\subsection{Boundary estimator}
$\Da^{*}_{\mathrm{disc}}$, and $R^{*}$ under H2a, is the 50\,\% crossing of an
unpenalised maximum-likelihood logistic regression of the replicate-level success outcomes
on $\log_{10}\Da$ (or $\log_{10}R$). The confidence interval is a percentile bootstrap
(2000 resamples) that resamples replicates \emph{within} each $\Da$ level. The design is
fixed, and only the stochastic replicates are resampled. No boundary is reported, and no
verdict word issued, when there is no crossing inside the tested range: all successes, all
failures, a fitted slope with the wrong sign, or a crossing off the grid. The tested range
is reported instead. A boundary is never extrapolated off the grid.
% p2/analysis.py boundary(), n_boot default 2000.

\textbf{Calibration, measured before the freeze.} On 100 simulated datasets with the
design's grid (9 $\Da$ levels, 20 replicates) and a known boundary, the estimator had
coverage 0.94 (nominal 0.95), bias $+0.012$ decade, and a standard deviation of the
estimate of 0.089 decade.
% PREREG_P2 s.4.4; tests/test_p2.py::test_boundary_ci_has_nominal_coverage_and_no_bias
% (docstring: "measured once at 100 datasets: coverage 0.94, bias +0.012 decade").

\subsection{Minimum detectable effect}
A boundary shift, or a gap between discovery and identifiability, is resolvable at about
$2.8 \times 0.089 \approx 0.25$ decade (80\,\% power, two-sided 5\,\%). \textbf{The MDE is
0.25 decade}, and it is quoted with every null. It applies on O1, where discovery and
identifiability both rest on 20 replicates per level. O2 and O3 identifiability has 3
replicates and is reported descriptively. The $\Da$ grid spacing is half a decade;
the logistic estimator resolves the crossing well inside one grid step, which is what the
measured standard deviation shows.
% PREREG_P2 s.4.4 and s.10 item 1 (the earlier "half a decade" was the grid spacing and is
% withdrawn).

\subsection{H2c refinement grid}
A factor-2 shift is 0.30 decade, close to the MDE. H2c is therefore tested on a refined
$\Da$ grid of \textbf{8 extra levels, log-uniform within $\pm 0.5$ decade of the best
non-EIV method's $\Da^{*}_{\mathrm{disc}}$} at O1, $\sigma = 2\,\%$, $\varepsilon = 2\,\%$,
for each isotherm. The 8 levels lie strictly inside $\pm 0.5$ decade; the endpoints
coincide with main-grid levels and are not repeated. These levels need new ground-truth
solves, made after the main grid and before any EIV method is run on them. The refinement
is part of the design, not a post-hoc rescue.
% PREREG_P2 s.4.4; s.4.3 "H2c resolution".

\subsection{How each hypothesis is decided}
\begin{description}
\item[H2a.] $R$ is computed per cell. Because $\sigma$ and $\varepsilon$ enter $R$ itself,
they are not slicing axes. A slice is one ($\Pe$, isotherm, observation model) combination,
pooling every $\sigma$, $\varepsilon$ and $\Da$ in it. $P(\text{success})$ of the best
method (\S\ref{sec:metrics}) is fitted against $\log_{10} R$ by the same logistic
estimator, and $R^{*}$ is each slice's 50\,\% crossing. Only O1 and O2 enter: on O3 the
sole discovery method is M1's single-replicate demonstration, which yields no $R^{*}$. The
spread is $\max - \min$ of $\log_{10} R^{*}$ over slices, with a bootstrap interval. Slices without a crossing are listed, never imputed. The
verdict words follow the thresholds in \S\ref{sec:hyp}. When the boundary does not
collapse, the axis that separates the $R^{*}$ values is named.
% PREREG_P2 s.4.3 "R and its collapse (H2a)"; p2/analysis.py r_collapse.
\item[H2b.] For each O1 slice ($\sigma$, $\varepsilon$, isotherm), the gap is
$\log_{10}\Da^{*}_{\mathrm{disc}} - \log_{10}\Da^{*}_{\mathrm{ident}}$. Here
$\Da^{*}_{\mathrm{disc}}$ comes from the best method, with its selection-carrying bootstrap
interval (\S\ref{sec:metrics}), and $\Da^{*}_{\mathrm{ident}}$ from the same logistic
estimator (\S\ref{sec:ident}). Gap $\le -0.25$ gives \verdict{discovery fails before
identifiability}; gap $\ge +0.25$ gives \verdict{discovery outlives identifiability};
otherwise \verdict{discovery fails with identifiability}. O2 and O3 gaps are reported
descriptively, with intervals and no verdict words.
% p2/analysis.py disc_vs_ident (mde 0.25); PREREG_P2 s.2 H2b.
\item[H2c.] H2c is decided at $\sigma = 2\,\%$ on O1. The shift is
\[
\log_{10}\Da^{*}_{\mathrm{disc}}(\text{M5}) - \log_{10}\Da^{*}_{\mathrm{disc}}(\text{best
strong-form non-EIV}),
\]
with a bootstrap interval. It is estimated on the refined grid for $\varepsilon = 2\,\%$
and on the main grid for $\varepsilon \in \{0.5\,\%, 5\,\%\}$. The shift relative to the
best non-EIV method over both forms is reported alongside, descriptively. The
hypothesis holds at a given $\varepsilon$ if the shift is at least $\log_{10}2 = 0.30$
decade. A shift at or above the MDE but below 0.30 decade takes the intermediate verdict
words of \S\ref{sec:hyp}. The hypothesis is falsified if the shift is below the MDE at every
$\varepsilon > 0$.
\item[H2d.] Exploratory. We report, as a rate per cell, in which ($\sigma$, observation)
cells the loading-dependent rate is discovered under the L2 criterion of
\S\ref{sec:metrics}, with no directional prediction and no boundary hypothesis.
\end{description}

\subsection{External anchors}
First, S1 on the map. The Damk\"ohler number of \citep{santana2023} will be recomputed
under the definition of \S\ref{sec:axes}, and also in the residence-time form, and plotted.
Our literature review estimated its residence-time value at about 0.9, which is in the
kinetically controlled range. If S1 lands in our failure region, the contradiction is
reported and diagnosed, not explained away.
% research/lit_equation_discovery.md S1 addendum: Da = kL/v ~ 0.86 (our computation).
Second, real MOF-303 data (exploratory). The best method per regime is applied to published
MOF-303 water-uptake kinetics \citep{abuelmaaty2024} and to the digitised breakthrough curve
the companion paper already uses. No hypothesis rests on either.

% ---------------------------------------------------------------------
\section{Pilots: what was measured and what it changed}
\label{sec:pilots}
Pilot data are excluded from every confirmatory analysis. No discovery method has been run
on any grid cell. The pilots fall into two kinds. Generator pilots measured the simulator.
Feasibility and verification pilots ran methods only on known-answer systems or on fresh
solves that are not grid cells.

\paragraph{P1, generator pilot.}
\emph{Measured:} for $k$ on a log grid, the residence-time and run Damk\"ohler numbers,
grid convergence ($N_z$ = 400 against 800), breakthrough and cost.
% results/p2_pilot_generator.json.
\emph{Changed:}
(i) The primary Damk\"ohler number is $k\,t_{\mathrm{stoich}}$, not $kL/u$.
(ii) At a horizon of $1.5\,t_{\mathrm{stoich}}$ the default column reached only
$c/c_{\mathrm{in}} \approx 0.75$ at the outlet. The cause is the thermal wave: the outlet
was still about 9\,K above feed. At $k = 0.01$\,s$^{-1}$ the 95\,\% crossing was at
$2.2\,t_{\mathrm{stoich}}$, and a $4\,t_{\mathrm{stoich}}$ horizon reached
$c/c_{\mathrm{in}} = 1.00$. The default horizon is therefore $4\,t_{\mathrm{stoich}}$,
subject to a per-solve check that the exit reaches 0.95; a solve that does not is
extended, never silently truncated.
(iii) The exit curve changed by at most 0.11\,\% between $N_z$ = 400 and 800 at every $k$
tested ($10^{-4}$--$10$\,s$^{-1}$), so $N_z = 400$ is adequate.
(iv) The run Damk\"ohler number $k\,t_{\mathrm{final}}$ spans about 1 to $10^{5}$ over that
$k$ range. The targeted $\Da$ range $10^{-1}$--$10^{3}$ therefore corresponds to $k
\approx 10^{-5}$--$10^{-1}$\,s$^{-1}$ at default physics.
(v) The residence-time Damk\"ohler number is $0.4\,\mathrm{s}\times k$, from $4\times10^{-5}$
to 4 over the pilot's $k$ range, which is why it is not the axis.
% PREREG_P2 s.6 items 1-4, 7 and s.4.2; results/p2_pilot_generator.json (exit_final
% 0.74-0.77 at 1.5 t_stoich; grid_change_rel <= 0.00112; Da_res 4e-5 ... 4).
The ground-truth solves have since been generated, at 3--7\,s each. All 81 reached
95\,\% breakthrough within their horizon, some after extension to $8\,t_{\mathrm{stoich}}$.
Generating ground truth runs no discovery method and is permitted before the freeze.
% p2_generate.log: "done: 81 cells; without 95 % breakthrough: none"; horizon 8 entries.

\paragraph{P2, M1 timing pilot.}
\emph{Measured} on one cell at $\Da \approx 14$: the differentiable column matches the
reference solver to below $5\times10^{-3}$ on the exit curve, and $N_z = 50$ differs from
$N_z = 400$ by 0.55\,\%. On the study CPU a forward solve takes 22\,s, a gradient with
respect to $k$ 101\,s, and a gradient with respect to a $2\times32$ network 145\,s. That
allows about 50 training steps in the 2\,h budget, where a UDE needs hundreds to thousands.
\emph{Changed:} by the rule declared before the pilot, M1 is not run at scale
(\S\ref{sec:compet}).
% results/p2_pilot_m1.json: Da 14.40, forward_s 21.5, grad_k_s 101.4, grad_nn_s 145.2,
% exit_change_50_vs_400 0.0055, epochs_in_2h 49.6.

\paragraph{P3, M1b verification.}
\emph{Measured} on a fresh solve that is not a grid cell: from dense exact $c$ ($N_z$ =
200), the inversion recovers $q$ to within 3\,\% of $q_{\max}$. The median error was 0.013
against a median driving force of 0.18, unchanged at $N_z = 800$, so the error is not
numerical diffusion. Fed the true $q$, the weak-form pipeline recovers the law exactly at
$\Da \approx 14$. Fed the inverted $q$, it does not; it selects $c^{2}$ and $T$. The
inversion error concentrates at the front, where the kinetic signal is.
\emph{Changed:} nothing in the methods or thresholds. The fact is declared because it bears
on H2a: for O2, $\varepsilon_{\mathrm{eff}}$ must include the state-reconstruction error,
and $\delta$ will be reported against it. The definition of $\varepsilon_{\mathrm{rec}}$ in
\S\ref{sec:hyp} implements this.
% PREREG_P2 s.6 item 5.

\paragraph{P4, sampling resolution of the O1 profile.}
\emph{Measured} on a fresh noise-free non-isothermal run: the profile's best $k$ at an
interior probe sits $+0.61\,\%$ high with linear interpolation between 200 snapshots,
$+0.47\,\%$ with monotone cubic interpolation, and $+0.17\,\%$ at 800 snapshots.
\emph{Changed:} monotone cubic interpolation was adopted. The bias is immaterial for the
$[k/2, 2k]$ criterion and is reported with the O1/O2 results. The discovery methods face
the same sampling.
% PREREG_P2 s.6 item 6.

\paragraph{P5, known-answer tests that forced design facts.}
These are M4's BIC residual-variance floor (and its recalibration to $2\times10^{-3}$ at
the grid's 200-sample resolution), the adaptive profile-likelihood refinement, M2's
admissibility rule of at least 3 equations per library term, the documented M8 failure, and
the ODR-BINDy accuracy gap (\S\ref{sec:compet}). Each was found on a known-answer system and
is declared before the freeze.

\paragraph{P7, M8 cost gate.}
\emph{Measured:} one M8 fit at the grid's size took 1144\,s, against the 300\,s limit
written in the pilot script before it ran. The machine was shared with other long jobs at
the time. \emph{Changed:} M8 is not run on the grid
(\S\ref{sec:compet}).
% results/p2_cost_gates.json.

\paragraph{P6, pre-freeze revisions.}
Drafting this manuscript against the pre-registration and the code exposed 20 places where
the two disagreed or where the pre-registration was silent. Each was fixed in the
pre-registration in the open, in a dated list of pre-freeze revisions, before any discovery
method ran on any grid cell. The fixes were of three kinds. Some aligned the text to
the implemented and tested harness: the library, unconstrained M4, the $z$-derivatives and
the adaptive refinement. Some defined quantities that had been named but not specified:
$\delta$, $\varepsilon_{\mathrm{eff}}$, $R$, the best method, $\Da^{*}_{\mathrm{ident}}$,
the success criteria for M9, L2 and Lib-A, and the observation model. Others added verdict
words for the gaps between thresholds, the placement rule for the H2c levels, and the M5
statistic. One withdrew a stale MDE statement. One tightened a test: M6 now checks every
coefficient to 5\,\%. One implemented M2's support sweep rather than dropping it.
A second pass the same day did six things. It left one definition of
$\Da^{*}_{\mathrm{ident}}$ and made H2b confirmatory on O1 only. It defined the H2a slices,
$\varepsilon_{\mathrm{rec}}$, the operational L2 success rule and the Lib-A fidelity. It
fixed H2c's resolution for each $\varepsilon$, gave the isothermal control a grid, and
replaced M5's grid method with EIV best-subset. It also added the O2 profile likelihood
and grid runners for every method. A third pass left one version each of the M5, M8 and
L2 rules. It made H2c like for like (strong form, O1), wrote out the libraries for M6 and
M9, and restricted the 0.25-decade MDE to O1.
% PREREG_P2 s.10 items 1-20, "Second pass", "Third pass"; p2_grid_o1.py, p2_grid_o23.py (o2, o2ident,
% o3), p2_grid_extra.py (m5, m6, m7, m8, m9, l2).

\paragraph{Pending before the freeze.}
Two items remain. The first is the ODR-BINDy five-seed check, which decides how the port is
reported on its own example: as verified, or as ``port, accuracy below the reference''. The
second is the M7 inclusion gate, which decides whether M7 runs or is reported NOT RUN. Both
rules were written before the runs that decide them. The pre-registration is frozen by a commit after
these and before any discovery method touches a grid cell. Any later change is a dated,
reasoned amendment, reported as post hoc if made after results exist.

% ---------------------------------------------------------------------
\section{Timeline and compute}
\label{sec:compute}
All computation runs on one CPU workstation without a GPU, one training job at a time,
through a job queue with a completion test per job. The confirmatory grid drivers refuse to
run until the freeze commit exists and the pre-registration is unchanged since it. There is
no override.
\begin{itemize}
\item Ground truth: the 81 solves and the 9 isothermal-control solves are complete; the
H2c refinement adds 8 $\Da$ levels per isotherm after the main grid.
\item O1 and O2 regression methods (M2--M6, M9, and M1b followed by the O1 methods) at 20
replicates over about 2600 cells. These take seconds each. Every method has a grid runner,
and every runner refuses to start before the freeze.
% p2_grid_o1.py, p2_grid_o23.py, p2_grid_extra.py.
% PREREG_P2 s.7. 2600 ~ 81 x 16 (sigma, eps) x 2 (O1, O2) = 2592.
\item M7, if admitted, on its O1 subgrid, and the one-replicate M1 demonstration: about 200
searches or trainings in total.
\item O2 and O3 identifiability: a full column solve per profile point, 3 replicates on
the subgrid.
\end{itemize}
Order of work: the remaining pre-freeze items (\S\ref{sec:pilots}), then the freeze, then
Stage-1 review. The grid runs after in-principle acceptance, or after the internal freeze if
the author decides not to wait; the Stage-2 report states which.
% RESEARCH_PROGRAM.md "Format decision (2026-10-04)".

% ---------------------------------------------------------------------
\section{Limitations}
\label{sec:limits}
\begin{itemize}
\item \emph{Synthetic data.} The ground truth is a one-dimensional plug-flow-with-dispersion
model. Non-plug flow, which can corrupt the extraction of LDF and dispersion coefficients
\citep{knox2016}, is outside the model. O3 is the realistic case and is reported first; O1
is a ceiling.
\item \emph{Competitor strength.} The claim is only as strong as the weakest
strongest-configuration method. M1 is run only as a demonstration because of compute. On
the grid, M5 is a classical errors-in-variables estimator, not a published EIV discovery
method. ODR-BINDy and WENDy are written for autonomous systems and are reported only on
their own examples. M7 is gated. M8 fails its known-answer test and is not run on the grid
by its cost gate, measured on a shared machine. Each is reported in those terms. Lib-B is smaller than first planned, so
the boundaries are optimistic with respect to library size. Two methods
recommended by our literature review are not in the competitor set: physics-constrained
weak-form symbolic regression \citep{servia2024,servia2025} and partitioning by time-scale
regime \citep{muhammed2025}. Their absence is a limitation, not a finding.
\item \emph{Axis choice.} $\Da$ may not be the right single axis \citep{hansen2017,
miyabe2000}. That is why $\Pe$ and isotherm shape are axes and why H2a is stated in terms
of $R$.
\item \emph{Sampling.} Stiffness relative to the sampling interval breaks discovery in
chemical kinetics \citep{prabhu2025}. The sampling here is fixed by the observation model,
so at the highest $\Da$ the kinetic time is short relative to the sampling interval. This
is reported, not varied.
\item \emph{Selection of the best method} per slice favours discovery, and so is
conservative for the verdict \verdict{discovery fails before identifiability}. Its
interval carries the selection (\S\ref{sec:metrics}).
\item \emph{Single machine.} The compute limits are those of one CPU. They decide which
methods run at scale, not what the methods can do.
\end{itemize}

% ---------------------------------------------------------------------
\section{Data and code availability}
The solver, harness (\texttt{p2/}), observation model, grid drivers, pilots, tests and the
pre-registration with its full version history will be released in the study repository.
Ground-truth fields are regenerated deterministically from the manifest. Every corrupted
observation is a deterministic function of (cell, replicate, $\sigma$, $\varepsilon$,
observation model), so any reported number can be reproduced from the code. Pilot results
files are released with the Stage-1 submission. Confirmatory results files will be released
with Stage 2.

\section*{Declaration of competing interest}
The author declares no competing interests.

\section*{Funding}
This research received no specific grant from
any funding agency in the public, commercial or not-for-profit sectors.

\section*{Declaration of generative AI and AI-assisted technologies}
During this work the author used Claude (Anthropic), an AI assistant, extensively: to
write and test the method harness, the observation model, the grid drivers and the
analysis code; to draft and revise the pre-registration and this Stage-1 manuscript; to
retrieve and check bibliographic records; and to act as an internal reviewer. The
pre-registration's version history and the test suite record what was built and when.
The author reviewed and edited the content and takes full responsibility for the
publication.

% ---------------------------------------------------------------------
\begin{thebibliography}{99}
\bibitem{brunton2016} Steven L. Brunton, Joshua L. Proctor, J. Nathan Kutz, Discovering governing equations from data by sparse identification of nonlinear dynamical systems, Proceedings of the National Academy of Sciences, vol. 113, 3932-3937 (2016). \href{https://doi.org/10.1073/pnas.1517384113}{doi:10.1073/pnas.1517384113}; \href{https://arxiv.org/abs/1509.03580}{arXiv:1509.03580}.

\bibitem{rudy2016} Samuel H. Rudy, Steven L. Brunton, Joshua L. Proctor, J. Nathan Kutz, Data-driven discovery of partial differential equations, arXiv preprint arXiv:1609.06401 (2016). \href{https://arxiv.org/abs/1609.06401}{arXiv:1609.06401}.

\bibitem{messenger2021ode} Daniel A. Messenger, David M. Bortz, Weak SINDy: Galerkin-Based Data-Driven Model Selection, Multiscale Modeling \& Simulation, vol. 19, 1474-1497 (2021). \href{https://doi.org/10.1137/20M1343166}{doi:10.1137/20M1343166}; \href{https://arxiv.org/abs/2005.04339}{arXiv:2005.04339}.

\bibitem{messenger2021pde} Daniel A. Messenger, David M. Bortz, Weak SINDy For Partial Differential Equations, arXiv preprint arXiv:2007.02848 (2020). \href{https://arxiv.org/abs/2007.02848}{arXiv:2007.02848}.

\bibitem{reinbold2020} Patrick A. K. Reinbold, Daniel R. Gurevich, Roman O. Grigoriev, Using Noisy or Incomplete Data to Discover Models of Spatiotemporal Dynamics, Phys. Rev. E 101, 010203 (2020). \href{https://arxiv.org/abs/1911.03365}{arXiv:1911.03365}.

\bibitem{fasel2022} U. Fasel, J. N. Kutz, B. W. Brunton, S. L. Brunton, Ensemble-SINDy: Robust sparse model discovery in the low-data, high-noise limit, with active learning and control, Proceedings of the Royal Society A: Mathematical, Physical and Engineering Sciences, vol. 478, 20210904 (2022). \href{https://doi.org/10.1098/rspa.2021.0904}{doi:10.1098/rspa.2021.0904}; \href{https://arxiv.org/abs/2111.10992}{arXiv:2111.10992}.

\bibitem{hirsh2021} Seth M. Hirsh, David A. Barajas-Solano, J. Nathan Kutz, Sparsifying Priors for Bayesian Uncertainty Quantification in Model Discovery, arXiv preprint arXiv:2107.02107 (2021). \href{https://arxiv.org/abs/2107.02107}{arXiv:2107.02107}.

\bibitem{kaheman2020} Kadierdan Kaheman, J. Nathan Kutz, Steven L. Brunton, SINDy-PI: a robust algorithm for parallel implicit sparse identification of nonlinear dynamics, Proceedings of the Royal Society A: Mathematical, Physical and Engineering Sciences, vol. 476, 20200279 (2020). \href{https://doi.org/10.1098/rspa.2020.0279}{doi:10.1098/rspa.2020.0279}; \href{https://arxiv.org/abs/2004.02322}{arXiv:2004.02322}.

\bibitem{bertsimas2022} Dimitris Bertsimas, Wes Gurnee, Learning Sparse Nonlinear Dynamics via Mixed-Integer Optimization, arXiv preprint arXiv:2206.00176 (2022). \href{https://arxiv.org/abs/2206.00176}{arXiv:2206.00176}.

\bibitem{cranmer2023} Miles Cranmer, Interpretable Machine Learning for Science with PySR and SymbolicRegression.jl, arXiv preprint arXiv:2305.01582 (2023). \href{https://arxiv.org/abs/2305.01582}{arXiv:2305.01582}.

\bibitem{udrescu2020} Silviu-Marian Udrescu, Max Tegmark, AI Feynman: a Physics-Inspired Method for Symbolic Regression, Sci. Adv. 6:eaay2631 (2020). \href{https://arxiv.org/abs/1905.11481}{arXiv:1905.11481}.

\bibitem{santana2023} Vinicius V. Santana, Erbet Costa, Carine M. Rebello, Ana Mafalda Ribeiro, Christopher Rackauckas, Idelfonso B.R. Nogueira, Efficient hybrid modeling and sorption model discovery for non-linear advection-diffusion-sorption systems: A systematic scientific machine learning approach, Chemical Engineering Science, vol. 282, 119223 (2023). \href{https://doi.org/10.1016/j.ces.2023.119223}{doi:10.1016/j.ces.2023.119223}; \href{https://arxiv.org/abs/2303.13555}{arXiv:2303.13555}.

\bibitem{valocchi1985} Albert J. Valocchi, Validity of the Local Equilibrium Assumption for Modeling Sorbing Solute Transport Through Homogeneous Soils, Water Resources Research, vol. 21, 808-820 (1985). \href{https://doi.org/10.1029/wr021i006p00808}{doi:10.1029/wr021i006p00808}.

\bibitem{hansen2017} Scott K. Hansen, Velimir V. Vesselinov, Local equilibrium and retardation revisited, arXiv preprint arXiv:1703.03087 (2017). \href{https://arxiv.org/abs/1703.03087}{arXiv:1703.03087}.

\bibitem{katsoulas2023} Konstantinos Katsoulas, Monica Tirapelle, Eva Sorensen, Luca Mazzei, On the apparent dispersion coefficient of the equilibrium dispersion model: An asymptotic analysis, Journal of Chromatography A, vol. 1708, 464345 (2023). \href{https://doi.org/10.1016/j.chroma.2023.464345}{doi:10.1016/j.chroma.2023.464345}.

\bibitem{knox2016} James C. Knox, Armin D. Ebner, M. Douglas LeVan, Robert F. Coker, James A. Ritter, Limitations of Breakthrough Curve Analysis in Fixed-Bed Adsorption, Industrial \& Engineering Chemistry Research, vol. 55, 4734-4748 (2016). \href{https://doi.org/10.1021/acs.iecr.6b00516}{doi:10.1021/acs.iecr.6b00516}.

\bibitem{ward2022} Adam Ward, Ronny Pini, Integrated uncertainty quantification and sensitivity analysis of single-component dynamic column breakthrough experiments, Adsorption, vol. 28, 161-183 (2022). \href{https://doi.org/10.1007/s10450-022-00361-z}{doi:10.1007/s10450-022-00361-z}.

\bibitem{taylor2026} W. Taylor, J.D. Herman, V.L. Morales, Global sensitivity of laboratory-scale groundwater reactive transport simulations for sorption and degradation: Breakthrough dynamics across Péclet--Damköhler regimes, Journal of Contaminant Hydrology, vol. 283, 105043 (2026). \href{https://doi.org/10.1016/j.jconhyd.2026.105043}{doi:10.1016/j.jconhyd.2026.105043}.

\bibitem{miyabe2000} Kanji Miyabe, Georges Guiochon, Determination of the lumped mass transfer rate coefficient by frontal analysis, Journal of Chromatography A, vol. 890, 211-223 (2000). \href{https://doi.org/10.1016/s0021-9673(00)00537-9}{doi:10.1016/s0021-9673(00)00537-9}.

\bibitem{stanhope2014} S. Stanhope, J. E. Rubin, D. Swigon, Identifiability of Linear and Linear-in-Parameters Dynamical Systems from a Single Trajectory, SIAM Journal on Applied Dynamical Systems, vol. 13, 1792-1815 (2014). \href{https://doi.org/10.1137/130937913}{doi:10.1137/130937913}.

\bibitem{gallo2026} Matteo Gallo, Fabio Anselmi, Paolo Lazzari, Attractor Geometry Determines the Identifiability Limits of System Discovery, arXiv preprint arXiv:2607.18490 (2026). \href{https://arxiv.org/abs/2607.18490}{arXiv:2607.18490}.

\bibitem{delgado2025} Diemen Delgado-Cano, Erick Kracht, Urban Fasel, Benjamin Herrmann, SINDy on slow manifolds, arXiv preprint arXiv:2507.00747 (2025). \href{https://arxiv.org/abs/2507.00747}{arXiv:2507.00747}.

\bibitem{klishin2024} Andrei A. Klishin, Joseph Bakarji, J. Nathan Kutz, Krithika Manohar, Statistical Mechanics of Dynamical System Identification, arXiv preprint arXiv:2403.01723 (2024). \href{https://arxiv.org/abs/2403.01723}{arXiv:2403.01723}.

\bibitem{muhammed2025} Ismaila Muhammed, Dimitris M. Manias, Dimitris A. Goussis, Haralampos Hatzikirou, Data-driven identification of biological systems using multi-scale analysis, PLOS Computational Biology, vol. 21, e1013193 (2025). \href{https://doi.org/10.1371/journal.pcbi.1013193}{doi:10.1371/journal.pcbi.1013193}.

\bibitem{prabhu2025} Siddharth Prabhu, Nick Kosir, Mayuresh V. Kothare, Srinivas Rangarajan, Derivative-Free Domain-Informed Data-Driven Discovery of Sparse Kinetic Models, Industrial \& Engineering Chemistry Research, vol. 64, 2601-2615 (2025). \href{https://doi.org/10.1021/acs.iecr.4c02981}{doi:10.1021/acs.iecr.4c02981}.

\bibitem{huang2021} Juntao Huang, Yizhou Zhou, Wen-An Yong, Data-driven discovery of multiscale chemical reactions governed by the law of mass action, Journal of Computational Physics, vol. 448, 110743 (2022). \href{https://doi.org/10.1016/j.jcp.2021.110743}{doi:10.1016/j.jcp.2021.110743}; \href{https://arxiv.org/abs/2101.06589}{arXiv:2101.06589}.

\bibitem{delahunt2021} Charles B. Delahunt, J. Nathan Kutz, A toolkit for data-driven discovery of governing equations in high-noise regimes, arXiv preprint arXiv:2111.04870 (2021). \href{https://arxiv.org/abs/2111.04870}{arXiv:2111.04870}.

\bibitem{kaptanoglu2023} Alan A. Kaptanoglu, Lanyue Zhang, Zachary G. Nicolaou, Urban Fasel, Steven L. Brunton, Benchmarking sparse system identification with low-dimensional chaos, arXiv preprint arXiv:2302.10787 (2023). \href{https://arxiv.org/abs/2302.10787}{arXiv:2302.10787}.

\bibitem{nogueira2022} Idelfonso B. R. Nogueira, Vinicius V. Santana, Ana M. Ribeiro, Alírio E. Rodrigues, Using scientific machine learning to develop universal differential equation for multicomponent adsorption separation systems, The Canadian Journal of Chemical Engineering, vol. 100, 2279-2290 (2022). \href{https://doi.org/10.1002/cjce.24495}{doi:10.1002/cjce.24495}.

\bibitem{frandsen2025} Jesper Frandsen, Vinicius V. Santana, Peter Jul-Rasmussen, Idelfonso B.R. Nogueira, Jakob K. Huusom, Krist V. Gernaey, Jens Abildskov, A systematic screening of neural network-based hybrid models of adsorption in chromatography processes, AIChE Journal, vol. 71, e70045 (2025). \href{https://doi.org/10.1002/aic.70045}{doi:10.1002/aic.70045}.

\bibitem{jessop2026} Arthur Jessop, Mohammed Alsubeihi, Ben Moseley, Ashwin Kumar Rajagopalan, Differentiable Hybrid Modelling for Learning and Optimising Chemical Transport Processes from Experimental Data, arXiv preprint arXiv:2609.04011 (2026). \href{https://arxiv.org/abs/2609.04011}{arXiv:2609.04011}.

\bibitem{praditia2022} Timothy Praditia, Matthias Karlbauer, Sebastian Otte, Sergey Oladyshkin, Martin V. Butz, Wolfgang Nowak, Learning Groundwater Contaminant Diffusion-Sorption Processes With a Finite Volume Neural Network, Water Resources Research, vol. 58, e2022WR033149 (2022). \href{https://doi.org/10.1029/2022WR033149}{doi:10.1029/2022WR033149}.

\bibitem{subraveti2022} Sai Gokul Subraveti, Zukui Li, Vinay Prasad, Arvind Rajendran, Can a computer ``learn'' nonlinear chromatography?: Physics-based deep neural networks for simulation and optimization of chromatographic processes, Journal of Chromatography A, vol. 1672, 463037 (2022). \href{https://doi.org/10.1016/j.chroma.2022.463037}{doi:10.1016/j.chroma.2022.463037}.

\bibitem{scheurer2026} Stefania Scheurer, Riccardo Frenner, Tim Brünnette, Sergey Oladyshkin, Wolfgang Nowak, Efficient confidence interval computation for physics-aware machine learning of diffusion-sorption models, Frontiers in Water, vol. 8, 1813791 (2026). \href{https://doi.org/10.3389/frwa.2026.1813791}{doi:10.3389/frwa.2026.1813791}.

\bibitem{haghpanah2023} Reza Haghpanah, Danny Shade, Fitting Adsorption Isotherms with Symbolic Regression, Industrial \& Engineering Chemistry Research, vol. 62, 22141-22148 (2023). \href{https://doi.org/10.1021/acs.iecr.3c02900}{doi:10.1021/acs.iecr.3c02900}.

\bibitem{sharlin2024} Samiha Sharlin, Tyler R. Josephson, In Context Learning and Reasoning for Symbolic Regression with Large Language Models, arXiv preprint arXiv:2410.17448 (2024). \href{https://arxiv.org/abs/2410.17448}{arXiv:2410.17448}.

\bibitem{cornelio2021} Cristina Cornelio, Sanjeeb Dash, Vernon Austel, Tyler Josephson, Joao Goncalves, Kenneth Clarkson, Nimrod Megiddo, Bachir El Khadir, Lior Horesh, AI Descartes: Combining Data and Theory for Derivable Scientific Discovery, arXiv preprint arXiv:2109.01634 (2021). \href{https://arxiv.org/abs/2109.01634}{arXiv:2109.01634}.

\bibitem{fox2023} Charles Fox, Neil Tran, Nikki Nacion, Samiha Sharlin, Tyler R. Josephson, Incorporating Background Knowledge in Symbolic Regression using a Computer Algebra System, arXiv preprint arXiv:2301.11919 (2023). \href{https://arxiv.org/abs/2301.11919}{arXiv:2301.11919}.

\bibitem{servia2024} Miguel Ángel de Carvalho Servia, Ilya Orson Sandoval, King Kuok (Mimi) Hii, Klaus Hellgardt, Dongda Zhang, Ehecatl Antonio del Rio Chanona, The automated discovery of kinetic rate models -- methodological frameworks, Digital Discovery, vol. 3, 954-968 (2024). \href{https://doi.org/10.1039/d3dd00212h}{doi:10.1039/d3dd00212h}; \href{https://arxiv.org/abs/2301.11356}{arXiv:2301.11356}.

\bibitem{servia2025} Miguel Ángel de Carvalho Servia, Ilya Orson Sandoval, King Kuok, Hii, Klaus Hellgardt, Dongda Zhang, Ehecatl Antonio del Rio Chanona, Constraint-Guided Symbolic Regression for Data-Efficient Kinetic Model Discovery, arXiv preprint arXiv:2507.02730 (2025). \href{https://arxiv.org/abs/2507.02730}{arXiv:2507.02730}.

\bibitem{loman2025} Torkel E Loman, Ruth E Baker, Functional and parametric identifiability for universal differential equations applied to chemical reaction networks, arXiv preprint arXiv:2510.14140 (2025). \href{https://arxiv.org/abs/2510.14140}{arXiv:2510.14140}.

\bibitem{katsoulas2026} Konstantinos Katsoulas, Federico Galvanin, Luca Mazzei, Eva Sorensen, A Diagnostic Procedure for Identifying Isotherm Models in Liquid Chromatography, Industrial \& Engineering Chemistry Research, vol. 65, 1277-1292 (2026). \href{https://doi.org/10.1021/acs.iecr.5c03704}{doi:10.1021/acs.iecr.5c03704}.

\bibitem{abuelmaaty2024} Ahmed E. Abu EL-Maaty, Mahmoud A. Abdalla, Mohamed Essalhi, Mahmoud M. Abdelnaby, Morsi M. Mahmoud, Mohamed A. Habib, Mohamed Antar, Rached Ben-Mansour, Kinetics of Water Adsorption in Metal-Organic Framework(MOF-303) for Adsorption Cooling Application, Energy Conversion and Management: X, vol. 24, 100694 (2024). \href{https://doi.org/10.1016/j.ecmx.2024.100694}{doi:10.1016/j.ecmx.2024.100694}.

\bibitem{coyle2026} Reid A. Coyle, Shyam Chand Pal, Peter Walther, Saeun Park, Bin Feng, Zhiling Zheng, Sustainable Metal-Organic Framework Water Harvesters in the Artificial Intelligence Era, arXiv preprint arXiv:2605.29179 (2026). \href{https://arxiv.org/abs/2605.29179}{arXiv:2605.29179}.

\bibitem{raue2009} A. Raue, C. Kreutz, T. Maiwald, J. Bachmann, M. Schilling, U. Klingmüller, J. Timmer, Structural and practical identifiability analysis of partially observed dynamical models by exploiting the profile likelihood, Bioinformatics, vol. 25, 1923-1929 (2009). \href{https://doi.org/10.1093/bioinformatics/btp358}{doi:10.1093/bioinformatics/btp358}.

\bibitem{fung2026} Lloyd Fung, Overcoming error-in-variable or stiffness problem in SINDy-like data-driven model discovery, ACM Transactions on AI for Science, 3831701 (2026). \href{https://doi.org/10.1145/3831701}{doi:10.1145/3831701}.

\bibitem{bortz2023} David M. Bortz, Daniel A. Messenger, Vanja Dukic, Direct Estimation of Parameters in ODE Models Using WENDy: Weak-form Estimation of Nonlinear Dynamics, arXiv preprint arXiv:2302.13271 (2023). \href{https://arxiv.org/abs/2302.13271}{arXiv:2302.13271}.

\bibitem{paper1} A.~Batayhi, What binds a surrogate: a pre-registered falsification ladder for transfer to unseen adsorbent parameter sets in metal--organic framework water-adsorption columns, companion manuscript (2026).

\bibitem{liu2024} Ziming Liu, Yixuan Wang, Sachin Vaidya, Fabian Ruehle, James Halverson, Marin Solja\v{c}i\'c, Thomas Y. Hou, Max Tegmark, KAN: Kolmogorov-Arnold Networks, ICLR 2025 (arXiv 2024). \href{https://arxiv.org/abs/2404.19756}{arXiv:2404.19756}.

\bibitem{slote2026} Kevin Slote, Jeremie Fish, Erik Bollt, KANDy: Kolmogorov-Arnold Networks and Dynamical System Discovery, arXiv preprint arXiv:2602.20413 (2026). \href{https://arxiv.org/abs/2602.20413}{arXiv:2602.20413}.

\bibitem{spotorno2026} Enzo Nicolas Spotorno, Josafat Leal Filho, Antonio Augusto Medeiros Frohlich, Empirical Stability Analysis of Kolmogorov-Arnold Networks in Hard-Constrained Recurrent Physics-Informed Discovery, arXiv preprint arXiv:2602.09988 (2026). \href{https://arxiv.org/abs/2602.09988}{arXiv:2602.09988}.
\bibitem{golub1980} Gene H. Golub, Charles F. Van Loan, An analysis of the total least squares problem, SIAM Journal on Numerical Analysis, vol. 17(6), 883-893 (1980). \href{https://doi.org/10.1137/0717073}{doi:10.1137/0717073}.

\bibitem{golub1987} G. H. Golub, Alan Hoffman, G. W. Stewart, A generalization of the Eckart-Young-Mirsky matrix approximation theorem, Linear Algebra and its Applications, vol. 88-89, 317-327 (1987). \href{https://doi.org/10.1016/0024-3795(87)90114-5}{doi:10.1016/0024-3795(87)90114-5}.
\end{thebibliography}

\end{document}
